\newcommand{\oratio}{\pars{Oratio.}

\noindent Grátiæ tuæ, quǽsumus, Dómine, supplícibus tuis tríbue largitátem, ut, mandáta tua, te operánte, sectántes, consolatiónem vitæ præséntis accípiant et futúra gáudia comprehéndant.

\pars{Pro pace in universo mundo.} \scriptura{Sir. 50, 25; 2 Esdr. 4, 20; \textbf{H416}}

\vspace{-4mm}

\antiphona{II D}{temporalia/ant-dapacemdomine.gtex}

\vfill

\noindent Deus, a quo sancta desidéria, recta consília et iusta sunt ópera: da servis tuis illam, quam mundus dare non potest, pacem; ut et corda nostra mandátis tuis dédita, et hóstium subláta formídine, témpora sint tua protectióne tranquílla.

\noindent Per Dóminum nostrum Iesum Christum, Fílium tuum, qui tecum vivit et regnat in unitáte Spíritus Sancti, Deus, per ómnia sǽcula sæculórum.

\noindent \Rbardot{} Amen.}
\newcommand{\matversus}{\noindent \Vbardot{} Fili mi, atténde ad sapiéntiam meam.

\noindent \Rbardot{} Et prudéntiæ meæ inclína aurem tuam.}
\newcommand{\lectioi}{\pars{Lectio I.} \scriptura{Iac. 1, 1-18}

\noindent De Epístula beáti Iacóbi apóstoli.

\noindent Iacóbus, Dei et Dómini Iesu Christi servus, duódecim tríbubus, quæ sunt in dispersióne, salútem.

\noindent Omne gáudium existimáte, fratres mei, cum in tentatiónibus váriis incidéritis, sciéntes quod probátio fídei vestræ patiéntiam operátur;

\noindent patiéntia autem opus perféctum hábeat, ut sitis perfécti et íntegri, in nullo deficiéntes.

\noindent Si quis autem vestrum índiget sapiéntia, póstulet a Deo, qui dat ómnibus affluénter et non impróperat, et dábitur ei.

\noindent Póstulet autem in fide nihil hǽsitans; qui enim hǽsitat, símilis est flúctui maris, qui a vento movétur et circumfértur.

\noindent Non ergo ǽstimet homo ille quod accípiat áliquid a Dómino, vir duplex ánimo, incónstans in ómnibus viis suis.

\noindent Gloriétur autem frater húmilis in exaltatióne sua, dives autem in humilitáte sua, quóniam sicut flos feni transíbit.

\noindent Exórtus est enim sol cum ardóre et arefécit fenum, et flos eius décidit, et decor vultus eius depériit; ita et dives in itinéribus suis marcéscet.

\noindent Beátus vir, qui suffert tentatiónem, quia, cum probátus fúerit, accípiet corónam vitæ, quam repromísit Deus diligéntibus se.

\noindent Nemo, cum tentátur, dicat: “A Deo tentor”; Deus enim non tentátur malis, ipse autem néminem tentat.

\noindent Unusquísque vero tentátur a concupiscéntia sua abstráctus et illéctus; dein concupiscéntia, cum concéperit, parit peccátum; peccátum vero, cum consummátum fúerit, génerat mortem.

\noindent Nolíte erráre, fratres mei dilectíssimi. Omne datum óptimum et omne donum perféctum de sursum est, descéndens a Patre lúminum, apud quem non est transmutátio nec vicissitúdinis obumbrátio.

\noindent Voluntárie génuit nos verbo veritátis, ut simus primítiæ quædam creatúræ eius.}
\newcommand{\responsoriumi}{\pars{Responsorium 1.} \scriptura{\Rbardot{} 2 Tim. 4, 8 \Vbardot{} ibid. 1, 12; \textbf{H286}}

\vspace{-5mm}

\responsorium{II}{temporalia/resp-repositaestmihicorona-CROCHU.gtex}{}}
\newcommand{\lectioii}{\pars{Lectio II.} \scriptura{Nn. 1-3 : PG 34, 854-855}

\noindent Ex Libro sancti Macárii abbátis De oratióne.

\noindent Caput omnis boni stúdii et vertex óperum bonórum assidúitas oratiónis est, per quam in aliórum possessiónem bonórum indúcimur, eo qui nos ad ea vocat, manum cooperatrícem porrigénte.

\noindent Etenim communicátio mýsticæ virtútis et connéxio dispositiónis sanctitátis in Deum ipsiúsque mentis per ineffábilem in Dóminum amórem, in oratióne iis qui digni reddúntur, obvéniunt: \emph \emph{Dedísti enim,} ait, \emph{lætítiam in corde meo,} et ipse dicit Dóminus : \emph{Regnum cælórum intra vos est.}

\noindent Nam intra nos esse regnum, quid áliud putémus significáre quam cæléstem lætítiam Spíritus Sancti animábus dignis sensibíliter immíssam esse?

\noindent Ut præstántius céteris est opus oratiónis, ita maióri opus habet sollicitúdine qui in id incúmbit, ne qua illi subrépat a malítia decéptio.

\noindent Nam quo quisque maióris adipiscéndi boni curam cóncipit, eo maióri diábolus conátu ei obstáre conténdit.

\noindent Quare multam eum adhibére necésse erit vigilántiam, ut fructus caritátis, humilitátis, simplicitátis, bonitátis, discretiónis e patiénti oratiónis cultúra ipsi pleni uberésque progérminent:

\noindent ipsiúsque perspícuum in divína progréssum profectúmque demónstrent: sicque álios ad paris æmulatiónem stúdii próvocent.}
\newcommand{\responsoriumii}{\pars{Responsorium 2.} \scriptura{\Rbardot{} Ps. 85, 12-13 \Vbardot{} ibid., 13; \textbf{H88}}

\vspace{-5mm}

\responsorium{IV}{temporalia/resp-confitebortibidomine-CROCHU.gtex}{}}
\newcommand{\lectioiii}{\pars{Lectio III.} \scriptura{Nn. 1-3 : PG 34, 854-855}

\noindent Sine intermissióne orándum et perseveránter insisténdum oratióni esse, et divínus Apóstolus, et ipsíus Dómini oráculum prǽcipit dicéntis: \emph{Quanto magis Deus fáciet vindíctam clamántium ad se die ac nocte.} Et illud : \emph{Vigiláte et oráte.}

\noindent Opórtet ígitur ubíque oráre : nec umquam conátum remíttere. Quanto autem capitálius qui perseveráre in oratióne decrévit, opus aggréssus est, tanto ipsum plus certáminis suscípere, tantóque magis id eum ad extrémum perdúcere opórtet, eo quod perserverántiæ oratiónis multa óbicit impediménta právitas, somnum, acédiam, pondus córporis, cogitatiónum aversiónem, mentis instabilitátem, languórem et réliqua malítiæ stúdia.

\noindent Prætérea afflictiónes et ipsórum malórum dǽmonum insúltus veheménter impugnántium nos, et omni vi obstántium, ne appropinquáre Christo váleat Deum post veritátem incessánter quærens ánima.}
\newcommand{\responsoriumiii}{\pars{Responsorium 3.} \scriptura{\Rbardot{} Ps. 100, 1-2 \Vbardot{} ibid., 2; \textbf{H89}}

\vspace{-5mm}

\responsorium{III}{temporalia/resp-misericordiametiudicium-CROCHU-cumdox.gtex}{}}
\include{hebdomadaviii}
\include{feriavi}
